% ======================================================================
% TeXFrog source for the game chain of Lemma~\ref{lem:fs-hints}.
%
% This is the single source of truth for the pseudocode.  It feeds
%   * the consolidated figure below (plain pdflatex, no extra tooling), and
%   * the interactive HTML viewer:
%         texfrog html serve gameops/gameops.tex
%     (run from report/, with ../.venv activated).
%
% Line numbers are written by hand with \lno so that a given number denotes
% the same line in every game of the chain; the proof cites them.
% ======================================================================

%%% Game roster.  Order fixes the meaning of ranges such as G3-G6.
\tfgames{gameops}{G0, G1, G2, G3, G4, G5, G6}
\tfgamename{gameops}{G0}{\bG_0}
\tfgamename{gameops}{G1}{\bG_1}
\tfgamename{gameops}{G2}{\bG_2}
\tfgamename{gameops}{G3}{\bG_3}
\tfgamename{gameops}{G4}{\bG_4}
\tfgamename{gameops}{G5}{\bG_5}
\tfgamename{gameops}{G6}{\bG_6}

\tfdescription{gameops}{G0}{The \EUFCMA game of $\SIG[\Sigma]$, with a lazily sampled random oracle.}
\tfdescription{gameops}{G1}{The signing oracle overwrites the random-oracle table instead of re-using a stored value; $\badprog$ records this.}
\tfdescription{gameops}{G2}{The honest transcripts are sampled eagerly, so that $\mathbf{Initialize}$ runs $\Real_\Sigma(1^\secp, \numOSign)$.}
\tfdescription{gameops}{G3}{The honest transcripts are replaced by hint-assisted simulated ones; from here on the witness is unused.}
\tfdescription{gameops}{G4}{A guess $\fpguess$ of the index of the random-oracle query underlying the forgery is drawn; it does not influence the output.}
\tfdescription{gameops}{G5}{When the guess is wrong, the verifier is run on the challenge at the guessed position.}
\tfdescription{gameops}{G6}{The guess and the challenge move to $\mathbf{Initialize}$; this is the simulation run by the impersonator $\bdv$.}

\tfmacrofile{notation.tex}
\tfpreamble{preamble.tex}

\tfcommentary{gameops}{G0}{commentary/G0.tex}
\tfcommentary{gameops}{G1}{commentary/G1.tex}
\tfcommentary{gameops}{G2}{commentary/G2.tex}
\tfcommentary{gameops}{G3}{commentary/G3.tex}
\tfcommentary{gameops}{G4}{commentary/G4.tex}
\tfcommentary{gameops}{G5}{commentary/G5.tex}
\tfcommentary{gameops}{G6}{commentary/G6.tex}

\begin{tfsource}{gameops}
\begin{pchstack}
%%%
\begin{pcvstack}[boxed]
\procedure{$\mathbf{Initialize}$}{
  \lno{1} (x, w) \gets \Gen_R(1^\secp) \\
  \lno{2} \Qset \coloneqq \emptyset; \; \hctr \coloneqq 0; \; \sctr \coloneqq 0 \\
  \tfonly{G2}{\lno{3} \pcfor i = 1 \until \numOSign \pcdo \\}
  \tfonly{G2}{\lno{4} \pcind (\com_i, \mathsf{st}_i) \gets \Prover_1(x, w); \; \chall_i \sample \cC \\}
  \tfonly{G2}{\lno{5} \pcind \rsp_i \gets \Prover_2(\mathsf{st}_i, \chall_i); \; \pi_i \coloneqq (\com_i, \chall_i, \rsp_i) \\}
  \tfonly{G3,G4,G5,G6}{\lno{3'} \pcfor i = 1 \until \numOSign \pcdo \\}
  \tfonly{G3,G4,G5,G6}{\lno{4'} \pcind \hint_i \gets \hintdistr_x; \; \pi_i \gets \Sim(x, \hint_i) \\}
  \tfonly{G6}{\lno{6} \fpguess \sample \set{1 \until \numRO + 1} \\}
  \tfonly{G6}{\lno{7} \chall^\circ \sample \cC \\}
  \lno{8} \pcreturn \pk \coloneqq x
}
\pcvspace
\procedure{$\RO(u)$}{
  \lno{10} \pcif \HT[u] = \bot \pcthen \\
  \lno{11} \pcind \hctr \coloneqq \hctr + 1; \; \QT[\hctr] \coloneqq u \\
  \tfonly{G0,G1,G2,G3,G4,G5}{\lno{12} \pcind \HT[u] \sample \cC \\}
  \tfonly{G6}{\lno{12'} \pcind \pcif \hctr = \fpguess \pcthen \HT[u] \coloneqq \chall^\circ \quad \pcelse \HT[u] \sample \cC \\}
  \lno{13} \pcreturn \HT[u]
}
\end{pcvstack}
%%%
\pchspace[0pt]
\begin{pcvstack}[boxed]
\procedure{$\OSign(\msg)$}{
  \lno{20} \sctr \coloneqq \sctr + 1; \; \Qset \coloneqq \Qset \cup \set{\msg} \\
  \tfonly{G0,G1}{\lno{21} (\com_\sctr, \mathsf{st}_\sctr) \gets \Prover_1(x, w) \\}
  \tfonly{G2,G3,G4,G5,G6}{\lno{21'} \pcparse \pi_\sctr \text{ as } (\com_\sctr, \chall_\sctr, \rsp_\sctr) \\}
  \lno{22} u \coloneqq \com_\sctr \concat x \concat \msg \\
  \tfonly{G0,G1}{\lno{23} \chall_\sctr \sample \cC \\}
  \tfonly{G0,G1}{\lno{24} \pcif \HT[u] \neq \bot \pcthen \\}
  \tfonly{G0,G1}{\lno{25} \pcind \badprog \coloneqq \btrue \\}
  \tfonly{G0}{\lno{26} \pcind \chall_\sctr \coloneqq \HT[u] \\}
  \lno{27} \HT[u] \coloneqq \chall_\sctr \\
  \tfonly{G0,G1}{\lno{28} \rsp_\sctr \gets \Prover_2(\mathsf{st}_\sctr, \chall_\sctr) \\}
  \lno{29} \pcreturn (\com_\sctr, \chall_\sctr, \rsp_\sctr)
}
\pcvspace
\procedure{$\mathbf{Finalize}(\msg^*, \sig^*)$}{
  \lno{30} \pcparse \sig^* \text{ as } (\com^*, \chall^*, \rsp^*) \\
  \tfonly{G4,G5}{\lno{31} \fpguess \sample \set{1 \until \numRO + 1} \\}
  \lno{32} \pcif \msg^* \in \Qset \pcthen \pcreturn 0 \\
  \lno{33} u^* \coloneqq \com^* \concat x \concat \msg^* \\
  \tfonly{G4,G5,G6}{\lno{34} \text{let } j \text{ be such that } \QT[j] = u^* \\}
  \lno{35} \pcif \chall^* \neq \HT[u^*] \pcthen \pcreturn 0 \\
  \tfonly{G4,G5,G6}{\lno{36} \pcif j \neq \fpguess \pcthen \\}
  \tfonly{G4,G5,G6}{\lno{37} \pcind \badev \coloneqq \btrue \\}
  \tfonly{G5,G6}{\lno{38} \pcind \chall^* \coloneqq \HT[\QT[\fpguess]] \\}
  \lno{39} \pcreturn \Ver_\Sigma(x, \com^*, \chall^*, \rsp^*)
}
\end{pcvstack}
%%%
\end{pchstack}
\end{tfsource}

% ----------------------------------------------------------------------
% The consolidated figure that appears in the paper.  Lines shared by all
% seven games are printed plainly; lines belonging to only some of them are
% annotated with the games they occur in.  Everything below is generated
% from the single source above, so the figure cannot drift out of sync with
% the per-game renders or with the HTML viewer.
% ----------------------------------------------------------------------
\begin{figure}[htb]
    \centering
    % Tighten the gutter between the two stacks and the gap after the line
    % numbers: at their defaults the figure is a few points wider than
    % \textwidth.  (\resizebox is not an option here -- cryptocode's stacks
    % contain vertical-mode material.)
    \renewcommand{\pclnspace}{0.4em}%
    \tfrenderfigure{gameops}{G0,G1,G2,G3}
    \caption{%
        The games $\printgame{gm0}$--$\printgame{gm3}$ of the proof of
        Lemma~\ref{lem:fs-hints}, as a single listing: a line carrying no
        annotation belongs to every game, and a line annotated
        $\cmt \printgame{gm4}, \printgame{gm5}$ occurs only in the games
        named. Primed numbers mark a line that \emph{replaces} the
        unprimed one of the same number in the games where it occurs. We
        write $\pi_i \coloneqq (\com_i, \chall_i, \rsp_i)$ for the $i$-th
        transcript, as $\Real_\Sigma$ and $\HintSim_\Sigma$ do.
        %
        Reading the chain off the figure: $\printgame{gm0}$ is the \EUFCMA
        game of $\SIG[\Sigma]$; $\printgame{gm1}$ deletes line $26$, which is
        what makes the pair identical-until-$\badprog$; $\printgame{gm2}$
        replaces the on-demand prover calls by the eagerly sampled
        transcripts of $\Real_\Sigma(1^\secp, \numOSign)$ and
        $\printgame{gm3}$ by those of
        $\HintSim_\Sigma(1^\secp, \numOSign, \Sim)$, after which the witness
        $w$ is no longer used. 
    }\label{fig:games1}
\end{figure}

\begin{figure}[htb]
    \centering
    % Tighten the gutter between the two stacks and the gap after the line
    % numbers: at their defaults the figure is a few points wider than
    % \textwidth.  (\resizebox is not an option here -- cryptocode's stacks
    % contain vertical-mode material.)
    \renewcommand{\pclnspace}{0.4em}%
    \tfrenderfigure{gameops}{G4,G5,G6}
    \caption{%
        $\printgame{gm4}$ adds the guess $\fpguess$
        and the flag $\badev$ without affecting the output;
        $\printgame{gm5}$ adds line $38$; and $\printgame{gm6}$ moves the
        choice of $\fpguess$ and of the challenge into
        $\mathbf{Initialize}$, planting the latter at the $\fpguess$-th
        fresh $\RO$ query.
        %
        If $\adv$ makes fewer than $\fpguess$ fresh $\RO$ queries then
        $\QT[\fpguess]$ is undefined at line $38$; we use the convention
        $\HT[\bot] = \bot$, which is immaterial since line $38$ is reached
        only after $\badev$ has been set and is identical in
        $\printgame{gm5}$ and $\printgame{gm6}$.%
    }\label{fig:games2}
\end{figure}

